\subsection{The remaining torus fibre on the point-three cyclic cover}

\begin{statement}{A.198}{Proved}
{An explicit generalized-Jacobian comparison and a norm-one frame obstruction}
Let \(C:y^2=x^5-54x\), let \(W\subset C\) be the cyclic-eight
cover in A.192, and put
\[
 P=(3,9),\quad L=\Q(\theta),\quad\theta^4=54,\qquad
 \mathcal G=\operatorname{Jac}(C,C\setminus W),\quad J=\operatorname{Jac}(C).
\tag{A.198.1}
\]
The generalized Abel map is based at \(P\); its compact component
is \(R\mapsto[R-P]\), not \([R-\infty]\).
There are actual tori and morphisms over \(\Q\)
\[
\begin{aligned}
 \mathcal T&=\G_m\times\operatorname{Res}_{L/\Q}\G_m,
 &\mathcal T'&=\operatorname{Res}_{L/\Q}\G_m\times\G_m,\\
 h(u,v)&=(v^2,uN_{L/\Q}v),
 &\Phi:\mathcal G&\longrightarrow J\times\mathcal T',
\end{aligned}
\tag{A.198.2}
\]
where \(\Phi\) is an isogeny of degree sixteen.  On the pointed
curve its two torus coordinates are the literal units
\[
 b=\frac{x-\theta}{3-\theta},\qquad c=\frac y9;
 \qquad a=\frac x3=\frac{c^2}{N_{L/\Q}b}.
\tag{A.198.3}
\]
Its restriction to the boundary torus is \(h\).
The kernel is canonically \(J[2]\), with the branch Weil-pairing
map specified below.  It induces an isomorphism on \(T_\ell\)
for every odd \(\ell\), and the exact sequence
\[
 0\longrightarrow T_2\mathcal G\xrightarrow{T_2\Phi}
 T_2J\oplus T_2\mathcal T'\xrightarrow{\partial}J[2]
 \longrightarrow0.
\tag{A.198.4}
\]

Write \(\widehat{k^\times}_2=\varprojlim_m k^\times/(k^\times)^{2^m}\)
additively.  In each of the four compact classes of A.192, choose the
point or unit-tangent reference on \(W\) displayed in (A.198.13).
The relative pro-two affine class has a unique torus coordinate
\[
 (u,v)\in\widehat{\Q^\times}_2\oplus\widehat{L^\times}_2.
\tag{A.198.5}
\]
Its relative unit coordinates satisfy exactly
\[
                   \beta=2v,\qquad\gamma=u+Nv.
\tag{A.198.6}
\]
Thus the global unit data in that compact fibre are exactly
\(\beta\in2\widehat{L^\times}_2\) and arbitrary
\(\gamma\in\widehat{\Q^\times}_2\).  Each such pair has two
affine section classes, differing by
\((u,v)=(0,\kappa_{2,L}(-1))\).  The exact conditions at the
dyadic place in (A.198.15) select one of these two classes whenever
they hold.  This classifies the stated affine pro-two fibre with its
dyadic point or cusp image; matching at every other place remains a
separate requirement.

In the fibre relative to \(P\), the actual map to the original
pro-two Kummer coordinates is
\[
 A-\kappa_2(3)=8u-2Nv,\qquad
 B-\kappa_2(-2)=-2Nv.
\tag{A.198.7}
\]
The individual point-three frame modulo eight is exactly
\(Nv\in4\widehat{\Q^\times}_2\), with the torsion correction
proved below.

There is an explicit nonzero norm-one class \(v\) in (A.198.5)
which vanishes at two.  Extending it by zero at the odd coefficient
primes, and taking \(u=0\), gives an actual section of the full
affine-cover abelian quotient \(E_U/\overline{[H,H]}\), where
\(H=\pi_1(W_{\overline\Q})\).  Its full original Kummer pair
is \((3,-2)\), and its localization at two is the point-three
image.  At 571 it is excluded by the actual finite quotient
\[
 E_U/(\overline{[H,H]}H^4),\qquad
 \Pi/(\overline{[H,H]}H^4)
       \simeq(\Z/4)^8\rtimes C_{32},
\tag{A.198.8}
\]
of order \(2^{21}\) and nilpotency class twelve.
It therefore is not an everywhere locally point-or-cusp-matching
section of this quotient.  All section classes in this statement
are first compared in the cover group using a selected sheet;
the final finite obstruction is invariant under downstairs
geometric conjugacy by the fixed-sheet argument in the proof.
\end{statement}

\begin{proof}
\emph{Boundary lattices and the actual unit isogeny.}
The boundary consists of the six Weierstrass points
\(Q_0=(0,0),Q_\infty=\infty\), and
\(Q_r=(r,0)\) for \(r^4=54\).
Its torus is
\((\G_m^2\times\operatorname{Res}_{L/\Q}\G_m)/\G_m\).
Dividing every component by the component at infinity identifies
it with \(\mathcal T\) in (A.198.2).
The sign convention sends prescribed boundary values of a rational
function \(f\), normalized to one at infinity, to the generalized
divisor class of \(\operatorname{div}(f)\).

Put \(D_r=[Q_r-Q_\infty]\), including \(r=0\).
The five classes generate \(J[2]\), and their only
\(\F_2\)-linear relation is \(\sum_rD_r=0\), by the literal
branch-divisor calculation in A.186.  Consequently the full
lattice of principal degree-zero boundary divisors is
\[
 \mathcal L=2\Z^5+\Z(1,1,1,1,1)\subset\Z^5,
                         \qquad [\Z^5:\mathcal L]=16.
\tag{A.198.9}
\]
Indeed \(\operatorname{div}(x-r)=2(Q_r-Q_\infty)\)
and \(\operatorname{div}(y)=\sum_rQ_r-5Q_\infty\).
Conversely every principal boundary divisor gives the zero
combination of the \(D_r\), so its parity vector is zero or
the all-one vector; these displayed functions realize both cases.

Evaluation of these functions on degree-zero divisors supported
on \(W\) defines characters of \(\mathcal G\).  If a divisor
is \(\operatorname{div}(f)\) with \(f=1\) on the boundary,
Weil reciprocity makes every such evaluation one, so the map
descends to the generalized Jacobian.  The normalized evaluations
on \([R-P]\) are exactly (A.198.3).  The curve equation proves
\(c^2=aNb\).  The torus with coordinates \((a,b,c)\) subject
to this equation is isomorphic to \(\mathcal T'\) via
\((b,c)\mapsto(c^2/Nb,b,c)\).
On the boundary torus reciprocity gives
\(a=u^2,b=v^2,c=uNv\), proving the literal map \(h\).

The proper quotient \(q:\mathcal G\to J\), together with these
characters, gives \(\Phi=(q,b,c)\).
Its kernel lies in \(\mathcal T\) and is
\[
 C=\{(u,v):v\in\operatorname{Res}_{L/\Q}\mu_2, u=Nv\}.
\tag{A.198.10}
\]
This has order sixteen.  The map is onto: over an algebraic
closure first lift the desired point of \(J\) to \(\mathcal G\),
then use the surjective torus map \(h\) to correct the two
unit coordinates.  The same construction on an etale cover gives
surjectivity as an algebraic-group morphism and its finite kernel.
It is therefore the asserted isogeny.
In the five-sign notation for \(C\), the signs have product one.
The canonical map
\(e\mapsto(e_2(e,D_r))_r\) identifies \(J[2]\) with this
even-sign module.  Indeed the pairing is one on \((D_r,D_r)\)
and minus one on distinct branch classes.  The resulting five
vectors have their unique sum relation and span the four-dimensional
even-sign space.  All maps are Galois equivariant.

The pointed generalized Abel map identifies \(H^{\rm ab}\)
with \(T\mathcal G\).  Here this is the actual comparison of
\cite[Proposition A.8(ii),(iii), pp.232--234]{BIR_MochizukiAlbanese2012}:
after a finite splitting field the boundary components are
geometrically irreducible, and the functorial comparison descends
with their permutation action.  The geometric group is free of
rank nine, so its torsion-free abelianization is its full
abelianization.  This also gives the exact boundary sequence
\[
 0\longrightarrow M_2=\Z_2(1)\oplus
       \operatorname{Ind}_{G_L}^{G_\Q}\Z_2(1)
 \longrightarrow T_2\mathcal G\longrightarrow T_2J
 \longrightarrow0.
\tag{A.198.11}
\]
Since \(\Pi/H\) is cyclic, \(\Pi''\subset\overline{[H,H]}\).
Thus every use of this module below is through an actual arithmetic
quotient of the full metabelian tripod group.

\emph{Tate maps and the finite framing ambiguity.}
There is a unique morphism \(\Psi:J\times\mathcal T'\to\mathcal G\)
defined by \(\Psi(\Phi(g))=[2]g\).  This is independent of the
choice of \(g\), since \(C\) is killed by two; etale local
inverse branches of \(\Phi\) give a morphism which glues.
It satisfies \(\Psi\Phi=[2]\) and \(\Phi\Psi=[2]\).
For \(a=(a_m)\in T_2(J\times\mathcal T')\), the map in
(A.198.4) is explicitly
\(\partial(a)=\Psi(a_1)\in C\), followed by the branch-pairing
identification.  Equivalently, if \(g_m\) lifts \(a_m\) under
\(\Phi\), it is \([2^m]g_m\); the value is independent of
the lift and of \(m\).  The standard inverse multiplication
diagram proves exactness, including the final surjectivity.
For odd coefficients the kernel has invertible order, giving
the stated isomorphism.  Thus the same exact sequence holds
for full Tate modules with final quotient \(C\).

Globally, and at every finite place, \((T_2J)^{G_k}=0\):
the rational torsion is finite, and locally a sufficiently small
formal subgroup is a torsion-free lattice of finite index.
The same holds for the Tate modules of the tori, since the
roots of unity in a number field or a finite local extension are
finite.  These observations apply to full Tate modules as well.
They are the direct invariant argument of A.167.
Consequently a pair of cohomology classes in
\(H^1(k,TJ\oplus T\mathcal T')\) lifts through \(T\Phi\)
exactly when its image under the literal map
\(H^1(k,\partial)\) vanishes.  Every nonempty fibre is a torsor
under \(C(k)\).  Globally \(C(\Q)=\mu_2(L)\) has order two.
At two \(L\otimes\Q_2\) is the Eisenstein quartic, so
\(C(\Q_2)\) also has order two and the restriction map is
an isomorphism.  This is the finite framing ambiguity; it is
not discarded when passing to the unit characters.

\emph{Unique residual torus coordinates and all four references.}
The point \(P\) splits the arithmetic affine abelian extension.
For a fixed compact pro-two cohomology class represented by one
of the four references, subtract that reference class.  Exactness
of (A.198.11) gives a class in \(H^1(\Q,M_2)\), and it is unique
because \((T_2J)^{G_\Q}=0\).  The same uniqueness holds at each
finite completion.  Explicitly, a cobounder in \(T_2\mathcal G\)
whose values lie in \(M_2\) has invariant image in \(T_2J\),
so already belongs to \(M_2\).
Kummer and Shapiro, with the actual induced-module coset maps,
identify
\[
 H^1_{\rm cts}(\Q,M_2)=
          \widehat{\Q^\times}_2\oplus\widehat{L^\times}_2.
\tag{A.198.12}
\]
The inverse-limit comparison uses the finite zero-cochain orbits
of roots of unity, as in A.169, and makes no compactness claim
about these global multiplicative completions.
The norm is the corestriction map under this identification.

Use the local parameter \(y\) at \(Q_0\), and
\(w=-x^2/y\) at \(Q_\infty\), with global unit tangent one.
The four reference values, expressed using the origin infinity
only in the compact column, are
\[
\begin{array}{c|c|c|c|c}
 Q&j_C(Q)&b_Q&c_Q&\text{inertia in }M_2\\ \hline
 (3,9)&D&1&1&0\\
 (3,-9)&-D&1&-1&0\\
 Q_0&T&-\theta/(3-\theta)&1/9&(1,0)\\
 Q_\infty&0&(3-\theta)^{-1}&-1/9&(-1,-\operatorname{res}1)
\end{array}
\tag{A.198.13}
\]
The compact classes relative to \(P\) are obtained by subtracting
\(D\).  At \(Q_0\), \(x=-y^2/54+\cdots\).
At infinity, \(x=w^{-2}(1+\cdots)\) and
\(y=-w^{-5}(1+\cdots)\), by A.192.
These expansions prove every constant and inertia vector in the
table.  A factor \(1+w h(w)\) has every root in the formal
characteristic-zero ring, so its tangential Kummer class is zero.

For relative unit classes
\(\beta=\kappa_{2,L}(b)-\kappa_{2,L}(b_Q)\) and
\(\gamma=\kappa_2(c)-\kappa_2(c_Q)\), restriction of \(h\)
to Tate modules gives (A.198.6), and also
\[
                    \kappa_2(a)-\kappa_2(a_Q)=2u.
\tag{A.198.14}
\]
Here \(a_Q=c_Q^2/Nb_Q\) is the corresponding leading coefficient.
Conversely choose any half \(v\) of \(\beta\) and put
\(u=\gamma-Nv\).  This constructs the actual affine pro-two
section class in the given compact fibre.  The only choices of
half differ by \(\kappa_{2,L}(-1)\), since \(L\) has a real
embedding and its two-primary roots of unity are \(\{1,-1\}\).
Its norm is one, so the two classes have exactly the difference
stated in the theorem.  This also exhibits the nonzero element
of the cohomological framing kernel explicitly.

Let \(\beta_2,\gamma_2\) denote localization at two and let
\(\operatorname{res}:\widehat{\Q_2^\times}_2\to
\widehat{(L\otimes\Q_2)^\times}_2\).
By A.192 and the injectivity of its global compact localization,
a locally point-or-cusp-matching class in this compact fibre
has precisely the proper centre \(Q\) at two.  The complete
unit conditions for that local image are
\[
\begin{array}{c|c|c}
 Q&\text{condition on unit data}&\text{selected torus condition}\\ \hline
 (3,\pm9)&\beta_2=0,\ \gamma_2=0&u_2=v_2=0\\
 Q_0&\beta_2=0&v_2=0\\
 Q_\infty&5\beta_2=2\operatorname{res}\gamma_2
                                      &v_2=\operatorname{res}u_2.
\end{array}
\tag{A.198.15}
\]
These are equivalences within the stated compact fibre.
For the first two rows, a half of zero is zero or the local
class of minus one; the two global choices interchange them.
For the last, \(5v_2-\operatorname{res}\gamma_2\) is killed
by two, and again the two choices interchange its two values.
Exactly one has \(v_2=\operatorname{res}(\gamma_2/5)\),
which gives \(u_2=\gamma_2/5\).  The cusp parameter is
\(-\gamma_2/5\).  In the \(Q_0\) row it is \(\gamma_2\).
The local class of minus one is nonzero in the full two-adic
multiplicative completion of a finite local field.
This proves both existence and uniqueness of the compatible
framing in every row.  The invariant argument above concerns
finite places; no vanishing of Tate invariants at the real
place is asserted, and any real component condition is retained
separately when a Selmer domain is used.

\emph{The map to the original Kummer pair, including torsion.}
In the fibre relative to \(P\), the original functions satisfy
\[
             \frac t3=\frac{a^4}{Nb},\qquad
             \frac{1-t}{-2}=\frac1{Nb}.
\tag{A.198.16}
\]
Substituting (A.198.6) and (A.198.14) proves (A.198.7).
The combined-character difference is \(8(u-Nv)\), as required
by the actual cyclic-eight cover.
The separate modulo-eight frame requires \(2Nv\in8\widehat{\Q^\times}_2\).
A priori cancellation gives
\(Nv\in4\widehat{\Q^\times}_2+\mu_2\); the possible minus
sign must be checked.  Put \(\pi=\theta^3/9\).
Then \(\pi^4=24\) and \(\theta=\pi^3/4\).
At three this is an Eisenstein tame quartic, with residue field
\(\F_3\), and \(N\pi=-24\) has unit residue one.
Every unit norm has fourth-power residue, also one.
Since the principal units of \(\Q_3\) are uniquely
two-divisible, every completed norm has zero \(\mu_2\) unit
component in \(\widehat{\Q_3^\times}_2\).
The element \(\kappa_2(-1)+4w\) has nonzero such component
for every \(w\).  It cannot be a norm.  This excludes the
extra sign and proves the exact modulo-eight criterion.

\emph{A nonzero norm-one torus class invisible at two.}
For \(a=1,3,5\), put
\[
 r_a=\frac{a-\theta}{a+\theta},\quad
 A_a=\sum_{m\geq0}\frac{54^m}{(4m+1)a^{4m+1}},\quad
 B_a=\sum_{m\geq0}\frac{54^m}{(4m+3)a^{4m+3}}.
\tag{A.198.17}
\]
All sums converge in \(\Z_2\), with the tail from \(m=N\)
contained in \(2^N\Z_2\).  Set \(c_a=(A_a,B_a)^t\), and
define the actual two-adic coefficients
\[
 (\lambda_1,\lambda_3,\lambda_5)
 =\frac18\bigl(\det(c_3,c_5),\det(c_5,c_1),\det(c_1,c_3)\bigr).
\tag{A.198.18}
\]
The exact truncated sums and the stated tail bound give
\[
      (\lambda_1,\lambda_3,\lambda_5)
                      \equiv(38,29,103)\pmod{256}.
\tag{A.198.19}
\]
In particular the three coefficients are integral and
\(\lambda_5\) is odd.  The determinant relation is an exact
identity in \(\Q_2^2\), not an identity inferred from the
finite residue computation.

The automorphism \(\theta\mapsto-\theta\) takes \(r_a\)
to \(r_a^{-1}\), so each has absolute norm one.
At two, \(v_\theta(r_a-1)=5\).  The logarithm is injective
on \(1+\theta^5\mathcal O_{L,2}\), since the exponential
and logarithm converge inversely at valuation greater than
\(v_\theta(2)=4\).  Its literal convergent expansion is
\[
                \log(r_a)=-2\theta A_a-2\theta^3B_a.
\tag{A.198.20}
\]
The class
\[
                 v=\sum_{a=1,3,5}\lambda_a\kappa_{2,L}(r_a)
                     \in\widehat{L^\times}_2
\tag{A.198.21}
\]
therefore has norm zero and localizes to zero at two.
The convergence of scalar powers in this principal-unit group
and the injectivity of its logarithm prove the latter assertion
in the actual completion, including torsion.

At 571 the polynomial factors as
\((x-5)(x+5)(x^2+25)\); the last factor is irreducible modulo
571.  All factors are simple.  Thus
\(L\otimes\Q_{571}=\Q_{571}\times\Q_{571}\times K_2\),
where \(K_2/\Q_{571}\) is unramified quadratic.
Since \(N(5-\theta)=571\), the valuations of \(r_5\)
in these three factors are \((1,-1,0)\).
The other two ratios are units there, their numerator norms
being \(-53\) and 27.  The valuations of (A.198.21) are
\((\lambda_5,-\lambda_5,0)\), so \(v\ne0\).

Extend \(v\) by zero at every odd coefficient prime, using the
actual Chinese-remainder isomorphism of completed Kummer groups
from A.97.  The torus cocycle \((0,v)\), included in
\(T\mathcal G=H^{\rm ab}\), gives a continuous arithmetic
section relative to the point \(P\).  Its norm is zero, so
(A.198.7), now coefficient by coefficient, gives exactly the
full original pair \((3,-2)\).  Its local torus class at two
is zero, hence its image there is the point \(P\).

\emph{The finite local obstruction at 571.}
For this section the finite \(b\)-character modulo four has
valuation vector
\[
                              (2,2,0)\quad\text{in }(\Z/4)^3.
\tag{A.198.22}
\]
Every denominator \(3-\theta\) is a unit at 571.
For an ordinary \(W(\Q_{571})\)-point with integral \(x\),
at most one of the two linear factors can have positive
valuation, and the quadratic factor is a unit.  For nonintegral
\(x\), all three valuations are equal.  Neither pattern gives
(A.198.22).  This argument includes the actual normalized
valuation on the unramified quadratic factor.

The same exclusion holds for cusps.  At \(Q_0\) all three
values are units.  At infinity all three exponents are the same
multiple of the tangent valuation.  The only rational quartic
boundary points at 571 are the lifts of \(5,-5\).  At either,
\(x-\theta=y^2/f'(\theta)+\cdots\) in its one matching
component, with \(f'(\theta)=4\theta^4\) a unit; the other
two components are units.  Thus at most one valuation can be
nonzero.  These are all local point and cusp cases, and none
has (A.198.22).

For the finite arithmetic quotient use
\(K=\overline{[H,H]}H^4\), with powers and commutators closed.
It is characteristic in \(H\) and hence normal in \(E_U\).
The \(b\)-character modulo four factors through it after the
literal restriction to \(G_L\), so the preceding local test
is a test in this actual finite quotient.
With the Nielsen basis \(v=x,u=yx^{-3}\), the nine Schreier
generators are \(u_j=v^juv^{-j}\) and \(z=v^8\).
After quotienting by \(K\) they commute and each has order four.
Consequently
\[
        x\longmapsto v,\qquad y\longmapsto u_0v^3
\tag{A.198.23}
\]
gives the isomorphism in (A.198.8): \(v\) has order 32 and
cyclically shifts the eight lamps.  The map is onto and is an
isomorphism on the open subgroup generated by the nine
Schreier classes, proving equality of kernels and the order.
In \((\Z/4)[v]/(v^8-1)\) one has the literal identities
\[
           (v-1)^{11}=2\sum_{j=0}^7v^j\ne0,
                     \qquad (v-1)^{12}=0.
\tag{A.198.24}
\]
The successive commutator groups are the successive augmentation
images of the lamp module, so this proves class twelve.

Finally any downstairs geometric conjugacy to a local point or
cusp transports the selected sheet to a fixed sheet of its
finite eighth-root fibre.  The actual cover decomposition group
then supplies a rational local point or rational boundary point
of \(W\), with the same finite unit class up to cover-geometric
conjugacy.  The latter leaves the Kummer character class unchanged.
The exhaustive test above therefore excludes every downstairs
local point or cusp image as well.  No all-place local matching
is asserted for the constructed section; its explicit failure
at 571 is the conclusion.

The exact certificate is
\path{code/point3_genus2_toric_certificate.py}, with retained
output and data.  It verifies the integral lattice and torus
matrices, branch signs, unit identities and tangent constants,
four log truncations with the proved tail precision, the literal
571 factorization and exhaustive valuation patterns, and the
actual finite quotient and augmentation identities.
\end{proof}
